\capitulo{6}{Conclusiones}

El proyecto culmina con el despliegue de una aplicación web accesible y de uso público que alcanza los objetivos planteados inicialmente. 

Este desarrollo contribuirá al ámbito sanitario, promoviendo una mejor educación sobre el cuidado de la piel y la prevención del \textit{melanoma}, el cáncer cutáneo más evitable, mediante recomendaciones fácilmente aplicables en la vida cotidiana. Su interfaz intuitiva garantiza que cualquier usuario, incluso sin conocimientos tecnológicos, pueda navegar y aprovechar sus funcionalidades.

Además, el código fuente y las distintas fases de elaboración quedan disponibles en un repositorio público, garantizando la reproducibilidad y sirviendo de base para futuras mejoras.

\section{Aspectos relevantes.}

Uno de los principales desafíos de este proyecto fue la obtención y homogeneización de datos ambientales fiables a nivel provincial, en particular los valores del índice de radiación ultravioleta (UV) y temperaturas en tiempo real.  Desde el punto de vista técnico, se logró desarrollar una aplicación web completamente funcional con el \textit{framework Shiny}, integrando múltiples fuentes de datos abiertos, cuyos datos fueron limpiados, procesados y adaptados para su visualización en mapas interactivos. La implementación de módulos específicos para el cálculo de riesgo, la recomendación personalizada y la consulta de históricos climáticos permitió abordar el problema desde una perspectiva integral en un contexto de salud pública.

El enfoque preventivo y educativo adoptado en el diseño de la plataforma resultó clave para maximizar su utilidad y accesibilidad. La inclusión de contenidos divulgativos, una interfaz intuitiva y mensajes personalizados adaptados al tipo de piel y ubicación del usuario contribuyen a la concienciación sobre los efectos de la exposición solar, alineándose con los principios de la salud pública.

A nivel competencial, este proyecto ha supuesto una oportunidad para consolidar conocimientos en programación con \textit{R}, diseño de aplicaciones web, análisis de datos y representación geoespacial. La capacidad de gestionar datos reales, resolver obstáculos técnicos y tomar decisiones de diseño orientadas al usuario refuerza el valor formativo de esta experiencia.

En conjunto, los resultados alcanzados demuestran la viabilidad técnica y el valor divulgativo de la plataforma \textit{UV-Dermis} como herramienta de apoyo a la prevención del \textit{melanoma} cutáneo, con potencial de mejora y escalabilidad en el futuro.
